\documentclass[11pt,a4paper]{article}
\usepackage[T1]{fontenc}
\usepackage{lmodern,amsmath,amssymb,amsthm,mathtools,microtype}
\usepackage[margin=27mm]{geometry}
\usepackage[colorlinks=true,linkcolor=blue,citecolor=blue,urlcolor=blue]{hyperref}
\newtheorem{theorem}{Theorem}[section]
\newtheorem{corollary}[theorem]{Corollary}
\newtheorem{remark}[theorem]{Remark}
\newcommand{\Q}{\mathbb Q}
\newcommand{\Z}{\mathbb Z}
\DeclareMathOperator{\rad}{rad}
\setlength{\emergencystretch}{2em}
\title{A Conditional Height Bound for a Sum of Even Powers}
\author{}\date{9 November 2026}
\begin{document}\maketitle
\begin{abstract}
The rational $abc$ conjecture implies a linear logarithmic height
bound for a rational pair in terms of its sum of fourth powers.
In particular, $h(u)+h(v)\le3h(u^4+v^4)+C$ for every rational
pair. The proof explicitly accounts for a common denominator and
for a common factor of the two numerators. A general even-power
version and the precise dependence on the $abc$ constant are given.
\end{abstract}

\section{Heights and the assumed inequality}

For $r=a/b\in\Q$, with $a,b$ relatively prime integers and $b>0$,
put $h(r)=\log\max\{|a|,b\}$. For a rational pair write
$(1:u:v)=(D:A':B')$ with $D>0$ and
$\gcd(D,A',B')=1$, and put
\[
             h_2(u,v)=\log\max\{D,|A'|,|B'|\}.
\]
In particular, $h(u),h(v)\le h_2(u,v)$. For a positive integer
$n$, let $\rad(n)$ be the product of its distinct prime divisors.

We use the following precise consequence of the rational $abc$
conjecture. For every $\varepsilon>0$ there is a constant
$K_\varepsilon\ge1$ such that
\begin{equation}\label{eq:abc}
 c\le K_\varepsilon\rad(abc)^{1+\varepsilon}
\end{equation}
whenever $a,b,c$ are positive relatively prime integers with
$a+b=c$. This is the standard formulation, for example
\cite[Section~1]{Granville}. The argument below applies
\eqref{eq:abc} only to the triple $(A^d,B^d,A^d+B^d)$ with
$d$ even and $A,B$ nonzero coprime integers.

\section{The height inequality}

\begin{theorem}\label{thm:even-power}
Let $d\ge4$ be even. Choose $\varepsilon>0$ such that
\[
       a=\frac{d}{1+\varepsilon}-3>0,
\]
and assume \eqref{eq:abc} with constant $K_\varepsilon\ge1$.
Then for all $u,v\in\Q$, putting $s=u^d+v^d$, one has
\begin{align}
 h_2(u,v)&\le\frac1a h(s)+C_{d,\varepsilon},\label{eq:pair}\\
 h(u)+h(v)&\le\frac2a h(s)+2C_{d,\varepsilon},\label{eq:sum}\\
 C_{d,\varepsilon}
 &=\frac{1+a}{da}\log2+
                    \frac{\log K_\varepsilon}{a(1+\varepsilon)}.
 \nonumber
\end{align}
The same constant works uniformly for all signs, denominators,
and common numerator factors.
\end{theorem}
\begin{proof}
If both $u$ and $v$ are zero, all three heights vanish.
If one coordinate is zero and the other is $w$, then
$h_2(u,v)=h(w)=h(s)/d$. Since $0<a<d$, both asserted
inequalities follow. Hence assume that $uv\ne0$.

Choose the primitive common-denominator representation above and
write
\[
 g=\gcd(|A'|,|B'|)>0,\qquad A=A'/g,\qquad B=B'/g.
\]
Then
\[
 u=\frac{gA}{D},\qquad v=\frac{gB}{D},\qquad
 \gcd(A,B)=1,\qquad \gcd(g,D)=1.
\]
The last equality follows from $\gcd(D,A',B')=1$.
Put
\[
 H=\max\{|A|,|B|\},\qquad N=A^d+B^d>0.
\]
Because $d$ is even,
\begin{equation}\label{eq:N-size}
                         H^d\le N\le2H^d.
\end{equation}
Let
\[
 k=\gcd(N,D^d),\qquad m_0=N/k,\qquad n=D^d/k.
\]
These positive integers satisfy $\gcd(m_0,n)=1$. Moreover
$\gcd(g,n)=1$, since $n$ divides $D^d$. It follows that
\begin{equation}\label{eq:reduced-output}
         s=\frac{g^dN}{D^d}=\frac{m}{n},\qquad
              m=g^d m_0,\qquad \gcd(m,n)=1.
\end{equation}
Thus if $M=\max\{m,n\}$, then $h(s)=\log M$ and
$1\le m_0\le m\le M$.

Every prime dividing $N$ either divides $D$ or survives in $m_0$.
Consequently
\begin{equation}\label{eq:radical}
                    \rad(N)\le Dm_0.
\end{equation}
It is important that this bound uses $m_0=m/g^d$, not $m$:
primes coming only from the common numerator factor $g$ need
not divide $N$.

Since $A$ and $B$ are coprime, the positive integers
$A^d,B^d,N$ are pairwise coprime. The assumed $abc$ inequality,
\eqref{eq:radical}, and $|AB|\le H^2$ give
\[
 H^d\le N
 \le K_\varepsilon\bigl(|AB|\rad(N)\bigr)^{1+\varepsilon}
 \le K_\varepsilon(H^2Dm_0)^{1+\varepsilon}.
\]
Therefore
\begin{equation}\label{eq:first-H}
 H^{d/(1+\varepsilon)-2}
               \le K_\varepsilon^{1/(1+\varepsilon)}Dm_0.
\end{equation}
The identity $D^d=nN/m_0$ and \eqref{eq:N-size} imply
\[
                    D\le2^{1/d}H(n/m_0)^{1/d}.
\]
Substitution into \eqref{eq:first-H} and division by $H$ yield
\begin{equation}\label{eq:H}
 H^a\le2^{1/d}K_\varepsilon^{1/(1+\varepsilon)}
                          m_0^{(d-1)/d}n^{1/d}.
\end{equation}

On the other hand, \eqref{eq:reduced-output} gives
$g=(m/m_0)^{1/d}$, and hence
\[
 \max\{gH,D\}\le2^{1/d}H(M/m_0)^{1/d}.
\]
The left-hand side is exactly the exponential of $h_2(u,v)$,
because $(D:gA:gB)$ is primitive. Taking logarithms and using
\eqref{eq:H}, we obtain
\begin{align*}
 h_2(u,v)
 &\le \log H+\frac1d(\log2+\log M-\log m_0)\\
 &\le C_{d,\varepsilon}
   +\frac{d-1-a}{da}\log m_0
   +\frac1{da}\log n+\frac1d\log M.
\end{align*}
All three coefficients in the last line are nonnegative:
$a<d-3<d-1$. Replacing $\log m_0$ and $\log n$ by
$\log M$ gives coefficient
\[
            \frac{d-1-a}{da}+\frac1{da}+\frac1d=\frac1a.
\]
This proves \eqref{eq:pair}. Since each individual coordinate
height is at most the pair height, \eqref{eq:sum} follows.
\end{proof}

\begin{corollary}\label{cor:quartic-three}
Assume rational $abc$. There is a constant $C_0\ge0$ such that
for every rational pair $(u,v)$,
\[
                   h(u)+h(v)\le3h(u^4+v^4)+C_0.
\]
If $K_{1/11}\ge1$ is an $abc$ constant at exponent $12/11$,
one can take
\[
            C_0=\frac54\log2+\frac{11}4\log K_{1/11}.
\]
For the positive rational square
$q=(u^4+v^4)^2$, with $(u,v)\ne(0,0)$, this becomes
\[
                         h(u)+h(v)\le\frac32h(q)+C_0.
\]
More generally, any fixed multiplicative constant strictly greater
than two is valid in front of $h(u^4+v^4)$, after changing the
additive constant.
\end{corollary}
\begin{proof}
Take $d=4$, $\varepsilon=1/11$, so that $a=2/3$.
For rational $s$, one has $h(s^2)=2h(s)$.
Finally $2/a=2(1+\varepsilon)/(1-3\varepsilon)$ tends to
two as $\varepsilon\to0^+$.
\end{proof}

\begin{corollary}
Assume rational $abc$. There is a constant $C_1$ such that
\[
                   h(u)+h(v)\le h(u^6+v^6)+C_1
                     \qquad(u,v\in\Q).
\]
One may take $C_1=\tfrac12\log2+
\tfrac56\log K_{1/5}$.
\end{corollary}
\begin{proof}
Take $d=6$, $\varepsilon=1/5$, so that $a=2$.
\end{proof}

\section{Unconditional scope}
The finite-fiber assertion for $(u,v)\mapsto u^4+v^4$ is
unconditional: each nonzero fiber has smooth projective completion
$U^4+V^4=sW^4$, a genus-three curve, and thus has finitely many
rational points by Faltings' theorem. The zero fiber consists only
of $(0,0)$. This proves the existence of some finite-valued height
bound, but supplies neither a linear bound nor an effective one.
The preceding linear estimates use \eqref{eq:abc} as an additional
arithmetic hypothesis.

\begin{thebibliography}{9}
\bibitem{Granville}
A.~Granville, Rational and integral points on quadratic twists
of a given hyperelliptic curve, \emph{International Mathematics
Research Notices} (2007), article rnm027.
\url{https://doi.org/10.1093/imrn/rnm027}; author's manuscript,
\url{https://dms.umontreal.ca/~andrew/PDF/hyperelliptics.pdf}.
\end{thebibliography}
\end{document}
